\documentclass{beamer}

\usepackage{graphicx}
\usepackage{amsmath}
\usepackage{setspace}
\usepackage{float} 
\usepackage{multirow}
\usepackage{array}
\usepackage{enumerate}
\usepackage{booktabs}
\usepackage{indentfirst} 
\usepackage[style=mla]{biblatex}
\usepackage{subfigure}
\usepackage{hyperref}
\usepackage{textpos}

\definecolor{Turquoise3}{RGB}{0, 134, 139}
\renewcommand{\emph}[1]{{\color{Turquoise3}\textsl{#1}}}
\newcommand{\C}{\mathbb{C}} \newcommand{\F}{\mathbb{F}} \newcommand{\R}{\mathbb{R}} \newcommand{\Q}{\mathbb{Q}}
\newcommand{\N}{\mathbb{N}}
\newcommand{\myseries}[2]{$#1_1,#1_2,\dots,#1_#2$}
\newcommand{\nullspace}{~\\[15pt]}

\usetheme{Madrid}
\setbeamertemplate{navigation symbols}{}

\addtobeamertemplate{frametitle}{}{
\begin{textblock*}{100mm}(0.85\textwidth,-1cm)
\includegraphics[height=1cm]{logo.png}
\end{textblock*}}

\definecolor{themecolor}{RGB}{25,25,112} 

\usecolortheme[named=themecolor]{structure}

\setbeamertemplate{items}[default]


\title{VV285 RC Part I}
\subtitle{\textbf{Elements of Linear Algebra}\\``Matrices are just linear maps!"}
\institute[UM-SJTU JI]{Univerity of Michigan-Shanghai Jiao Tong University Joint Institute}
\author{Xingjian Zhang}

\begin{document}

\begin{frame}
    \titlepage
    \begin{center}
        \includegraphics[height=2cm]{logo2.png}
    \end{center}
\end{frame}

\begin{frame}
    \
    \frametitle{Overview of Linear Algebra}
    \begin{enumerate}
        \item Systems of Linear Equations
        \item Finite-Dimensional Vector Spaces
        \item Inner Product Spaces
        \item Linear Maps
        \item Matrices
        \item Theory of Systems of Linear Equations
        \item Determinants
    \end{enumerate}
\end{frame}

\section{Systems of Linear Equations}
\begin{frame}
    \frametitle{Overview}
    \begin{enumerate}
        \item Linear System
              \\ Homogeneous vs. Inhomogeneous
              \\ Underdetermined vs. Overdetermined
        \item Equivalency of Linear System
        \item The Gauß – Jordan Algorithm
        \item Diagonalizable (Existence and Uniqueness of Linear System)
        \item \textbf{Fundamental Lemma for Homogeneous Equations}
    \end{enumerate}
\end{frame}

\subsection{Linear System}
\begin{frame}
    \frametitle{Definition}

    A \emph{linear system} of $m$ (algebraic) equations in $n$ unknowns \myseries{x}{n}$\in V$ is a set of equations
    \begin{equation}
        \begin{aligned}
            a_{11}x_1+a_{12}x_2+\cdots+a_{1n}x_n & =b_1   \\
            a_{21}x_1+a_{22}x_2+\cdots+a_{2n}x_n & =b_2   \\
                                                 & \vdots \\
            a_{m1}x_1+a_{m2}x_2+\cdots+a_{mn}x_n & =b_m
        \end{aligned}
        \label{eq1}
    \end{equation}

    where \myseries{b}{m}$\in V$ and $a_{ij}\in\F,i=1,\ldots,m,~j=1,\ldots,n.$\\
    If $b_1=b_2=\cdots=b_m=0$, then \eqref{eq1} is called a \emph{homogeneous system}. Otherwise, it is called an \emph{inhomogeneous system.}

    If $m<n$ we say that the system in \emph{underdetermined}, if $m>n$ the system is called \emph{overdetermined}.
    A solution of a linear system of equations \eqref{eq1} is a tuple of elements (\myseries{y}{n})$\in V^n$ such that the predicate \eqref{eq1} becomes a true statement.\\[3pt]
\end{frame}

\subsection{Equivalence of Linear System}
\begin{frame}
    \frametitle{Definition}
    We say that two systems of linear equations are \emph{equivalent} if any solution
    of the first system is also a solution of the second system and vice-versa.
    Thus the systems
    \begin{equation*}
        \begin{aligned}
            x_1+3x_2-x_3    & =1 \\
            ~~~~~~-5x_2+x_3 & =1 \\
            ~~~~~~10x_2+x_3 & =1
        \end{aligned}
        \qquad\qquad\text{and}\qquad\qquad
        \begin{aligned}
            x_1 & =2 \\x_2&=0\\x_3&=1
        \end{aligned}
    \end{equation*} are \emph{equivalent}.
\end{frame}

\subsection{Gau\ss-Jordan Algorithm}
\begin{frame}
    \frametitle{Algorithm}
    The goal of the \emph{Gau\ss-Jordan algorithm} (also called Gau\ss ian elimination) is to transform a system
    \begin{equation*}
        \begin{tabular}{ccc|c}
            $\ast$ & $\ast$ & $\ast$ & $\diamond$ \\
            $\ast$ & $\ast$ & $\ast$ & $\diamond$ \\
            $\ast$ & $\ast$ & $\ast$ & $\diamond$ \\
        \end{tabular}
        \qquad\qquad\qquad\qquad
        \ast\in~\R~or~\C,~~\diamond\in V
    \end{equation*}

    first into the form \\
    \begin{equation}
        \begin{array}{c}
            \begin{tabular}{|ccc|c}
                1 & $\ast$ & $\ast$ & $\diamond$ \\ \cline{1-1}
            \end{tabular} \\
            \begin{tabular}{c|cc|c}
                0 & $1$ & $\ast$ & $\diamond$ \\ \cline{2-2}
            \end{tabular} \\
            \begin{tabular}{cc|c|c}
                0 & $0$ & $1$ & $\diamond$ \\
                \cline{3-3}
            \end{tabular}
        \end{array}
    \end{equation}

    and subsequently into \\
    \begin{equation}
        \begin{array}{c}
            \begin{tabular}{|ccc|c}
                1 & $0$ & $0$ & $\diamond$ \\ \cline{1-1}
            \end{tabular} \\
            \begin{tabular}{c|cc|c}
                0 & $1$ & $0$ & $\diamond$ \\ \cline{2-2}
            \end{tabular} \\
            \begin{tabular}{cc|c|c}
                0 & $0$ & $1$ & $\diamond$ \\
                \cline{3-3}
            \end{tabular}
        \end{array}
    \end{equation}
\end{frame}

\begin{frame}
    \frametitle{Elementary Row Manipulations}
    \textbf{Include:}
    \begin{enumerate}
        \item Swapping (interchanging) two rows,
        \item Multiplying each element in a row with a number,
        \item Adding a multiple of one row to another row.
    \end{enumerate}
    ~\\[15pt]

    \textbf{Result:} Transform a system into a equivalent system.
    Since each row represents an equation, we are essentially \textbf{manipulating equations}.
\end{frame}

\subsection{Existence and Uniqueness of Solutions}
\begin{frame}
    \frametitle{Diagonalizable}
    A system of m equations with n unknowns will have a unique solution if and only if it is \emph{diagonalizable}. i.e. It can be transformed into diagonal form.\\[15pt]
    \textbf{Remark:} \emph{Diagonalization} turns out to be an important topic in VV286, especially in terms of \emph{ordinary differential equation systems}.
\end{frame}

\subsection{Fundamental Lemma for Homogeneous Equations}
\begin{frame}
    \frametitle{Lemma}
    The homogeneous system
    \begin{align*}
        a_{11}x_1+a_{12}x_2+\cdots+a_{1n}x_n & =0       \\
                                             & ~~\vdots \\
        a_{m1}x_1+a_{m2}x_2+\cdots+a_{mn}x_n & =0
    \end{align*}
    of $m$ equations in $n$ real or complex unknowns \myseries{x}{n} has a \textbf{non-trivial} solution if $n>m$.\\[15pt]

    \textbf{Remark:} This fundamental lemma contributes to prove that any basis of a vector space has the same length.

\end{frame}

\section{Finite-Dimensional Vector Spaces}

\begin{frame}
    \frametitle{Overview}
    \begin{enumerate}
        \item Linear Independence
        \item Span
        \item Basis
        \item Dimension
        \item Basis Extension Theorem
        \item Sum of Vector Space
    \end{enumerate}
\end{frame}

\subsection{Linear Independence}
\begin{frame}
    \frametitle{Definition}
    Let $V$ be a real or complex vector space and \myseries{v}{n}$\in V$. Then the vectors \myseries{v}{n} are said to be \emph{independent} if for all \myseries{\lambda}{n}$\in\F$ \vspace*{-4mm}
    \[\sum_{k=1}^{n}\lambda_kv_k=0\qquad\Rightarrow\qquad\lambda_1=\lambda_2=\cdots=\lambda_n=0.\]

    A finite set $M\subset V$ is called an \emph{independent set} if the elements of $M$ are independent.
\end{frame}

\subsection{Span}
\begin{frame}
    \frametitle{Definition}
    Let \myseries{v}{n}$\in V$ and \myseries{\lambda}{n}$\in\F$. Then the expression
    \begin{equation*}
        \sum_{k=1}^{n}\lambda_kv_k=\lambda_1v_1+\cdots+\lambda_nv_n
    \end{equation*}
    is called a \emph{linear combination} of the vectors \myseries{v}{n}.\\
    The set
    \[\text{span}\{v_1,\ldots,v_n\}=\left\{y\in V:y=\sum_{k=1}^{n}\lambda_kv_k,\lambda_1,\ldots,\lambda_n\in\F\right\}\]
    is called the \emph{(linear) span} or the \emph{linear hull} of the vectors \myseries{v}{n}.

\end{frame}

\begin{frame}
    \frametitle{Lemma}
    The vectors \myseries{v}{n}$\in V$ are independent if and only if
    \textbf{none of them is contained in the span of all the others}.
    \\[15pt]
    (How to prove?)
\end{frame}

\subsection{Basis}
\begin{frame}
    \frametitle{Definition}
    Let $V$ be a real or complex vector space. An $n$-tuple $\mathcal{B}=(b_1,\ldots,b_n)\in V^n$ is called an \emph{(ordered and finite) basis} of $V$
    if every vector $v$ has a \textbf{unique} representation
    \[v=\sum_{i=1}^{n}\lambda_ib_i,\qquad\qquad\qquad\qquad\lambda_i\in\F.\]
    The numbers $\lambda_i$ are called the \emph{coordinates} of $v$ with respect to $\mathcal{B}$.
    \\[15pt]
    The tuple of vectors (\myseries{e}{n}), $e_i\in\R^n$,
    \begin{equation*}
        e_i=(0,\ldots 0,\mathop{1}_{\begin{subarray}{c}
                \uparrow\\  i\text{th}\\ \text{entry}\end{subarray}},0,\ldots,0),\qquad\qquad i=1,\ldots,n,
    \end{equation*}
    is called the \emph{standard basis} or \emph{canonical basis} of $\R^n$.
\end{frame}

\begin{frame}
    \frametitle{Theorem}
    Let $V$ be a real or complex vector space.\\ An $n$-tuple $\mathcal{B}=(b_1,\ldots,b_n)\in V^n$ is a basis of $V$ if and only if
    \begin{enumerate}
        \item the vectors \myseries{b}{n} are linearly independent, i.e., $\mathcal{B}$ is an independent set,
        \item $V=\text{span}~\mathcal{B}.$
    \end{enumerate}
    (How to prove?)
    \nullspace
    \textbf{Remark:} This theorem is more practical than the definition of basis when proving some set is a basis of some vector space.
\end{frame}

\subsection{Dimension}
\begin{frame}
    \frametitle{Definition}
    Let $V$ be a real or complex finite-dimensional vector space, $V\neq\{0\}$. \textbf{Then any basis of $V$ has the same length }(number of elements).
    \nullspace
    \textbf{Remark:} This theorem can be proved by contradiction (Use the definition of basis and the fundamental lemma for homogeneous equations). With such a limitation, we can then define the \emph{dimension} of vector space.
    \nullspace
    Let $V$ be a real or complex vector space. Then $V$ is called \emph{finite-dimensional} if either
    \begin{itemize}
        \item $V={0}$ or
        \item $V$ possesses a finite basis.
    \end{itemize}
    If $V$ is not finite-dimensional, we say that it is \emph{infinite-dimensional}.
\end{frame}



\end{document}